\documentclass[a4paper]{article}

\usepackage{graphicx}
\usepackage{amsmath,amssymb}
\usepackage{minted}
\usepackage{multirow}
\usepackage{float}
\usepackage{algorithm}
\usepackage{algpseudocode}
\usepackage{todonotes}
\usepackage{forest}
\usepackage{xstring}
\usepackage{xcolor}
\usepackage[most]{tcolorbox}
\usepackage{xparse}
\usepackage{natbib}

\usetikzlibrary{graphs,graphdrawing}
\usegdlibrary{layered}

% for external graphviz
\input{/home/donkarlo/Dropbox/repo/nd_language_ai_project/src/nd_language_ai/formal/computer/programming/tex/latex/code_clips/external_graphviz.tex}

% for tags
\input{/home/donkarlo/Dropbox/repo/nd_language_ai_project/src/nd_language_ai/formal/computer/programming/tex/latex/part/control_sequence/kind/semtag/semtag.tex}

% for links
\input{/home/donkarlo/Dropbox/repo/nd_language_ai_project/src/nd_language_ai/formal/computer/programming/tex/latex/code_clips/seehere.tex}

\title{Atuoregression}
\date{}
\author{Mohammad Rahmani}

\begin{document}
   \maketitle

   Autoregression is a method in which the current or future value of a sequence is predicted from its previous values. It models the temporal dependency between successive observations. An autoregressive model describes what is likely to occur next given the past.

   \bibliographystyle{plainnat}
   \bibliography{/home/donkarlo/Dropbox/repo/nd_shared_research/bibliography/bibliography.bib}
\end{document}
